%============================================================
% MTH 112 Project - Graphs of other trig functions
% Updated 202504
%============================================================

\section{Graphs of Other Trigonometric Functions} \label{section-graphs-of-other-trig-functions}
In this section, we will discuss the graphs of tangent, cosecant, and secant and their domains.

\textbook{\href{https://openstax.org/books/algebra-and-trigonometry-2e/pages/8-2-graphs-of-the-other-trigonometric-functions}{8.2}}

\subsection*{Preparation Exercises} \label{preparation-exercises-section-graphs-of-other-trig-functions}

Answer the following without using a calculator. These questions are intended to check the prerequisite skills needed to complete the rest of the lab.

\begin{myPrep}
	Evaluate $\tan\left(\frac{2\pi}{3}\right)$ using the definition of tangent: $\tan(\theta)=\frac{\sin(\theta)}{\cos(\theta)}$.
\end{myPrep}
\vfill

\begin{myPrep}
	Evaluate $\csc\left(-\frac{5\pi}{6}\right)$ using the definition of cosecant: $\csc(\theta)=\frac{1}{\sin(\theta)}$.
\end{myPrep}
\vfill

\begin{myPrep}
	Which of the six trigonometric functions (sine, cosine, tangent, cosecant, secant, and cotangent) are odd functions?
\end{myPrep}
\vfill

\begin{myPrep}
	What is the domain of the sine function? What is the domain of the cosine functions?
\end{myPrep}
\vfill

\begin{myPrep}
	Which values of $\theta$ make it so that $\sin(\theta)$ is $0$? Hint: there are infinitely many answers.
\end{myPrep}
\vfill

%======================================================
\newpage
%======================================================

\subsection*{Practice Exercises} \label{practice-exercises-section-graphs-of-other-trig-functions}



\begin{minipage}{0.9\textwidth}
		Recall that for a right triangle with with angle $\theta$, the \say{other trigonometric functions} can be found by the following definitions. Visit \href{http://tiny.cc/six-trig-functions}{this interactive Desmos graph} to see the six trigonometric lengths in action.
	\end{minipage}
	\begin{minipage}{0.1\textwidth}
		\hfill \qrcode[height=1cm]{http://tiny.cc/six-trig-functions}
	\end{minipage}


	\begin{itemize}
		\begin{multicols}{4}
			\item $\tan(\theta)=\frac{\sin(\theta)}{\cos(\theta)}$
			\item $\cot(\theta)=\frac{\cos(\theta)}{\sin(\theta)}$
			\item $\sec(\theta)=\frac{1}{\cos(\theta)}$
			\item $\csc(\theta)=\frac{1}{\sin(\theta)}$
		\end{multicols}
	\end{itemize}

\begin{myPractice}
	\begin{enumerate}
		\item The cosecant function, $\csc(\theta)=\frac{1}{\sin(\theta)}$, would be undefined when $\sin(\theta)=0$. When is $\sin(\theta)=0$? There are infinitely many answers: write your answers in a full sentence describing the pattern.
		\vfill
		\item Write your solutions to $\sin(\theta)=0$ in set-builder notation.
		\vfill
		\item Based on your answers to the above questions, what is the domain of the cosecant function? Write your answer in a sentence explaining the pattern.
		\vfill
	\end{enumerate}
\end{myPractice}

\begin{myPractice}
	\begin{enumerate}
		\item The secant function, $\sec(\theta)=\frac{1}{\cos(\theta)}$, would be undefined when $\cos(\theta)=0$. When is $\cos(\theta)=0$? There are infinitely many answers: write your answers in a full sentence describing the pattern.
		\vfill
		\item Write your solutions to $\cos(\theta)=0$ in set-builder notation.
		\vfill
		\item Based on your answers to the above questions, what is the domain of the secant function? Write your answer in a sentence explaining the pattern.
		\vfill
	\end{enumerate}
\end{myPractice}

%======================================================
\newpage
%======================================================

\begin{myPractice}
	\begin{enumerate}
		\item The tangent function, $\tan(\theta)=\frac{\sin(\theta)}{\cos(\theta)}$, would be undefined when $\cos(\theta)=0$. This is the same problem that the secant function had. So, what is the domain of the tangent function? Write your answer in a sentence explaining the pattern.
		\vfill
		\item The cotangent function, $\cot(\theta)=\frac{\cos(\theta)}{\sin(\theta)}$, would be undefined when $\sin(\theta)=0$. This is the same problem that the cosecant function had. So, what is the domain of the cotangent function? Write your answer in a sentence explaining the pattern.
		\vfill
	\end{enumerate}
\end{myPractice}

\begin{myPractice}
	Match the \say{other trigonometric} functions with their graphs, shown in Figures \ref{fig:csc}, \ref{fig:tan}, \ref{fig:sec}, and \ref{fig:cot}. Use the definitions of the functions and their domains to decide which goes with which. Check your answers with your favorite graphing program.
	\begin{multicols}{4}
		\begin{enumerate}
			\item $y=\tan(x)$
			\item $y=\cot(x)$
			\item $y=\sec(x)$
			\item $y=\csc(x)$
		\end{enumerate}
	\end{multicols}
	\vfill
	\begin{multicols}{4}
				\begin{center}\captionsetup{type=figure}\captionof{figure}{Graph A}
					\begin{tikzpicture}[scale=0.7]
 						\begin{axis}[
 							framed,
 							xmin=-7,xmax=7,
 							ymin=-7,ymax=7,
						    xtick={-6,-5,...,6},
						    ytick={-6,-5,...,6},
						    grid=both,
						    clip=false,
 							]
							\addplot[first,line width=1pt,samples=100,<->]expression[domain=-7:-6.427]{cosec(180/pi*x)};
							\addplot[first,line width=1pt,samples=100,<->]expression[domain=-6.14:-3.285]{cosec(180/pi*x)};
							\addplot[first,line width=1pt,samples=100,<->]expression[domain=-2.998:-0.143]{cosec(180/pi*x)};
							\addplot[first,line width=1pt,samples=100,<->]expression[domain=0.143:2.998]{cosec(180/pi*x)};
							\addplot[first,line width=1pt,samples=100,<->]expression[domain=3.285:6.14]{cosec(180/pi*x)};
							\addplot[first,line width=1pt,samples=100,<->]expression[domain=6.427:7]{cosec(180/pi*x)};
							\addplot[variable=\t,domain=-7:7,<->,dashed,line width=0.8pt,samples=10,color=third]({-2*pi},{t});
							\addplot[variable=\t,domain=-7:7,<->,dashed,line width=0.8pt,samples=10,color=third]({-pi},{t});
							\addplot[variable=\t,domain=-7:7,<->,dashed,line width=0.8pt,samples=10,color=third]({0},{t});
							\addplot[variable=\t,domain=-7:7,<->,dashed,line width=0.8pt,samples=10,color=third]({pi},{t});
							\addplot[variable=\t,domain=-7:7,<->,dashed,line width=0.8pt,samples=10,color=third]({2*pi},{t});
							\node[left,color=third,rotate=90] at (axis cs:-2*pi,-7) {$x=-2\pi$};
							\node[left,color=third,rotate=90] at (axis cs:-pi,-7) {$x=-\pi$};
							\node[left,color=third,rotate=90] at (axis cs:0,-7) {$x=0$};
							\node[left,color=third,rotate=90] at (axis cs:pi,-7) {$x=\pi$};
							\node[left,color=third,rotate=90] at (axis cs:2*pi,-7) {$x=2\pi$};
 						\end{axis}
					\end{tikzpicture}\label{fig:csc}
				\end{center}
				\begin{center}\captionsetup{type=figure}\captionof{figure}{Graph B}
					\begin{tikzpicture}[scale=0.7]
 						\begin{axis}[
 							framed,
							xmin=-7,xmax=7,
 							ymin=-7,ymax=7,
						    xtick={-6,-5,...,6},
						    ytick={-6,-5,...,6},
						    grid=both,
						    clip=false,
 							]
							\addplot[first,line width=1pt,samples=100,<->]expression[domain=-7:-4.854]{tan(180/pi*x)};
							\addplot[first,line width=1pt,samples=100,<->]expression[domain=-4.57:-1.713]{tan(180/pi*x)};
							\addplot[first,line width=1pt,samples=100,<->]expression[domain=-1.429:1.429]{tan(180/pi*x)};
							\addplot[first,line width=1pt,samples=100,<->]expression[domain=1.713:4.57]{tan(180/pi*x)};
							\addplot[first,line width=1pt,samples=100,<->]expression[domain=4.854:7]{tan(180/pi*x)};
							\addplot[variable=\t,domain=-7:7,<->,dashed,line width=0.8pt,samples=10,color=third]({-3*pi/2},{t});
							\addplot[variable=\t,domain=-7:7,<->,dashed,line width=0.8pt,samples=10,color=third]({-pi/2},{t});
							\addplot[variable=\t,domain=-7:7,<->,dashed,line width=0.8pt,samples=10,color=third]({pi/2},{t});
							\addplot[variable=\t,domain=-7:7,<->,dashed,line width=0.8pt,samples=10,color=third]({3*pi/2},{t});
							\node[left,color=third,rotate=90] at (axis cs:-3*pi/2,-7) {$x=-\frac{3\pi}{2}$};
							\node[left,color=third,rotate=90] at (axis cs:-pi/2,-7) {$x=-\frac{\pi}{2}$};
							\node[left,color=third,rotate=90] at (axis cs:pi/2,-7) {$x=\frac{\pi}{2}$};
							\node[left,color=third,rotate=90] at (axis cs:3*pi/2,-7) {$x=\frac{3\pi}{2}$};
 						\end{axis}
					\end{tikzpicture}\label{fig:tan}
				\end{center}
				\begin{center}\captionsetup{type=figure}\captionof{figure}{Graph C}
					\begin{tikzpicture}[scale=0.7]
 						\begin{axis}[
 							framed,
							xmin=-7,xmax=7,
 							ymin=-7,ymax=7,
						    xtick={-6,-5,...,6},
						    ytick={-6,-5,...,6},
						    grid=both,
						    clip=false,
 							]
							\addplot[first,line width=1pt,samples=100,<->]expression[domain=-7:-4.854]{sec(180/pi*x)};
							\addplot[first,line width=1pt,samples=100,<->]expression[domain=-4.57:-1.713]{sec(180/pi*x)};
							\addplot[first,line width=1pt,samples=100,<->]expression[domain=-1.429:1.429]{sec(180/pi*x)};
							\addplot[first,line width=1pt,samples=100,<->]expression[domain=1.713:4.57]{sec(180/pi*x)};
							\addplot[first,line width=1pt,samples=100,<->]expression[domain=4.854:7]{sec(180/pi*x)};
							\addplot[variable=\t,domain=-7:7,<->,dashed,line width=0.8pt,samples=10,color=third]({-3*pi/2},{t});
							\addplot[variable=\t,domain=-7:7,<->,dashed,line width=0.8pt,samples=10,color=third]({-pi/2},{t});
							\addplot[variable=\t,domain=-7:7,<->,dashed,line width=0.8pt,samples=10,color=third]({pi/2},{t});
							\addplot[variable=\t,domain=-7:7,<->,dashed,line width=0.8pt,samples=10,color=third]({3*pi/2},{t});
							\node[left,color=third,rotate=90] at (axis cs:-3*pi/2,-7) {$x=-\frac{3\pi}{2}$};
							\node[left,color=third,rotate=90] at (axis cs:-pi/2,-7) {$x=-\frac{\pi}{2}$};
							\node[left,color=third,rotate=90] at (axis cs:pi/2,-7) {$x=\frac{\pi}{2}$};
							\node[left,color=third,rotate=90] at (axis cs:3*pi/2,-7) {$x=\frac{3\pi}{2}$};
 						\end{axis}
					\end{tikzpicture}\label{fig:sec}
				\end{center}
				\begin{center}\captionsetup{type=figure}\captionof{figure}{Graph D}
					\begin{tikzpicture}[scale=0.7]
 						\begin{axis}[
 							framed,
							xmin=-7,xmax=7,
 							ymin=-7,ymax=7,
						    xtick={-6,-5,...,6},
						    ytick={-6,-5,...,6},
						    grid=both,
						    clip=false,
 							]
							\addplot[first,line width=1pt,samples=100,<->]expression[domain=-7:-6.427]{cot(180/pi*x)};
							\addplot[first,line width=1pt,samples=100,<->]expression[domain=-6.14:-3.285]{cot(180/pi*x)};
							\addplot[first,line width=1pt,samples=100,<->]expression[domain=-2.998:-0.143]{cot(180/pi*x)};
							\addplot[first,line width=1pt,samples=100,<->]expression[domain=0.143:2.998]{cot(180/pi*x)};
							\addplot[first,line width=1pt,samples=100,<->]expression[domain=3.285:6.14]{cot(180/pi*x)};
							\addplot[first,line width=1pt,samples=100,<->]expression[domain=6.427:7]{cot(180/pi*x)};
							\addplot[variable=\t,domain=-7:7,<->,dashed,line width=0.8pt,samples=10,color=third]({-2*pi},{t});
							\addplot[variable=\t,domain=-7:7,<->,dashed,line width=0.8pt,samples=10,color=third]({-pi},{t});
							\addplot[variable=\t,domain=-7:7,<->,dashed,line width=0.8pt,samples=10,color=third]({0},{t});
							\addplot[variable=\t,domain=-7:7,<->,dashed,line width=0.8pt,samples=10,color=third]({pi},{t});
							\addplot[variable=\t,domain=-7:7,<->,dashed,line width=0.8pt,samples=10,color=third]({2*pi},{t});
							\node[left,color=third,rotate=90] at (axis cs:-2*pi,-7) {$x=-2\pi$};
							\node[left,color=third,rotate=90] at (axis cs:-pi,-7) {$x=-\pi$};
							\node[left,color=third,rotate=90] at (axis cs:0,-7) {$x=0$};
							\node[left,color=third,rotate=90] at (axis cs:pi,-7) {$x=\pi$};
							\node[left,color=third,rotate=90] at (axis cs:2*pi,-7) {$x=2\pi$};
 						\end{axis}
					\end{tikzpicture}\label{fig:cot}
				\end{center}
	\end{multicols}
\end{myPractice}

% \subsection*{Definitions}\label{definitions-section-graphs-of-other-trig} %ISSUE: no definitions?
%======================================================
\newpage
%======================================================

\subsection*{Exit Exercises} \label{exit-exercises-section-graphs-of-other-trig-functions}

\begin{myExit}
Which trigonometric functions have a domain that is the set of all real numbers except integer multiples of $\pi$?
\end{myExit}
\vfill

\begin{myExit}
Which trigonometric function is defined to be $\frac{\cos(\theta)}{\sin(\theta)}$?
\end{myExit}
\vfill

\begin{myExit}
Which trigonometric functions have a domain that is the set of all real numbers except odd-integer multiples of $\frac{\pi}{2}$?
\end{myExit}
\vfill

\begin{myExit}
Shown in Figure \ref{fig:tan-lines} is a graph of a transformation of $y=\tan(x)$. What values of $C$ and $D$ create this transformation?
	\begin{center}\captionsetup{type=figure}\captionof{figure}{A transformation of $y=\tan(x)$}
		\begin{tikzpicture}
 			\begin{axis}[
 				framed,
				xmin=-7,xmax=7,
 				ymin=-7,ymax=7,
			    xtick={-6,-5,...,6},
			    ytick={-6,-5,...,6},
			    grid=both,
			    clip=false,
 				]
				\addplot[first,line width=1pt,samples=100,<->]expression[domain=-6.958:-4.124]{tan(180/pi*(x-pi/4))+2};
				\addplot[first,line width=1pt,samples=100,<->]expression[domain=-3.816:-0.983]{tan(180/pi*(x-pi/4))+2};
				\addplot[first,line width=1pt,samples=100,<->]expression[domain=-0.675:2.159]{tan(180/pi*(x-pi/4))+2};
				\addplot[first,line width=1pt,samples=100,<->]expression[domain=2.467:5.3]{tan(180/pi*(x-pi/4))+2};
				\addplot[first,line width=1pt,samples=100,<->]expression[domain=5.608:7]{tan(180/pi*(x-pi/4))+2};
				\addplot[variable=\t,domain=-7:7,<->,dashed,line width=0.8pt,samples=10,color=third]({-5*pi/4},{t});
				\addplot[variable=\t,domain=-7:7,<->,dashed,line width=0.8pt,samples=10,color=third]({-pi/4},{t});
				\addplot[variable=\t,domain=-7:7,<->,dashed,line width=0.8pt,samples=10,color=third]({3*pi/4},{t});
				\addplot[variable=\t,domain=-7:7,<->,dashed,line width=0.8pt,samples=10,color=third]({7*pi/4},{t});
				\node[left,color=third,rotate=90] at (axis cs:-5*pi/4,-7) {$x=-\frac{5\pi}{4}$};
				\node[left,color=third,rotate=90] at (axis cs:-pi/4,-7) {$x=-\frac{\pi}{4}$};
				\node[left,color=third,rotate=90] at (axis cs:3*pi/4,-7) {$x=\frac{3\pi}{4}$};
				\node[left,color=third,rotate=90] at (axis cs:7*pi/4,-7) {$x=\frac{7\pi}{4}$};
 			\end{axis}
		\end{tikzpicture}\label{fig:tan-lines}
	\end{center}
\end{myExit}
\exitlikert{the graphs of other trigonometric functions}